\documentclass[10pt,a4paper]{amsart}
\usepackage{amsmath,amssymb,amsthm,graphicx}
% Retained minimal preamble for the Stacks Project chapter more-algebra.tex (More on Algebra).
% The chapter file pinned-original-source.tex carries only \input{preamble} on its first line, so no
% preamble text lives inside the chapter itself.  The Stacks Project uses one single shared preamble for
% every chapter; the snapshot stages it identically as sources/stacks-chow/preamble.tex and
% sources/stacks-etale-cohomology/preamble.tex.  The block below is copied mechanically, line for line,
% from sources/stacks-chow/preamble.tex lines 14-18 (the Xy-pic setup) and 154-217 (the native theorem
% declarations together with the mathematical macros).  Nothing was added, removed or redefined; the
% wrapper class and the amsmath/amssymb/amsthm/graphicx base packages are supplied by the assembler.
% Chapter: more-algebra.tex -- More on Algebra.
% For commutative diagrams we use Xy-pic
\usepackage[all]{xy}

% We use 2cell for 2-commutative diagrams.
\xyoption{2cell}

% Theorem environments.
%
\theoremstyle{plain}
\newtheorem{theorem}[subsection]{Theorem}
\newtheorem{proposition}[subsection]{Proposition}
\newtheorem{lemma}[subsection]{Lemma}

\theoremstyle{definition}
\newtheorem{definition}[subsection]{Definition}
\newtheorem{example}[subsection]{Example}
\newtheorem{exercise}[subsection]{Exercise}
\newtheorem{situation}[subsection]{Situation}

\theoremstyle{remark}
\newtheorem{remark}[subsection]{Remark}
\newtheorem{remarks}[subsection]{Remarks}

\numberwithin{equation}{subsection}

% Macros
%
\def\lim{\mathop{\mathrm{lim}}\nolimits}
\def\colim{\mathop{\mathrm{colim}}\nolimits}
\def\Spec{\mathop{\mathrm{Spec}}}
\def\Hom{\mathop{\mathrm{Hom}}\nolimits}
\def\Ext{\mathop{\mathrm{Ext}}\nolimits}
\def\SheafHom{\mathop{\mathcal{H}\!\mathit{om}}\nolimits}
\def\SheafExt{\mathop{\mathcal{E}\!\mathit{xt}}\nolimits}
\def\Sch{\mathit{Sch}}
\def\Mor{\mathop{\mathrm{Mor}}\nolimits}
\def\Ob{\mathop{\mathrm{Ob}}\nolimits}
\def\Sh{\mathop{\mathit{Sh}}\nolimits}
\def\NL{\mathop{N\!L}\nolimits}
\def\CH{\mathop{\mathrm{CH}}\nolimits}
\def\proetale{{pro\text{-}\acute{e}tale}}
\def\etale{{\acute{e}tale}}
\def\QCoh{\mathit{QCoh}}
\def\Ker{\mathop{\mathrm{Ker}}}
\def\Im{\mathop{\mathrm{Im}}}
\def\Coker{\mathop{\mathrm{Coker}}}
\def\Coim{\mathop{\mathrm{Coim}}}

% Boxtimes
%
\DeclareMathSymbol{\boxtimes}{\mathbin}{AMSa}{"02}

%
% Macros for moduli stacks/spaces
%
\def\QCohstack{\mathcal{QC}\!\mathit{oh}}
\def\Cohstack{\mathcal{C}\!\mathit{oh}}
\def\Spacesstack{\mathcal{S}\!\mathit{paces}}
\def\Quotfunctor{\mathrm{Quot}}
\def\Hilbfunctor{\mathrm{Hilb}}
\def\Curvesstack{\mathcal{C}\!\mathit{urves}}
\def\Polarizedstack{\mathcal{P}\!\mathit{olarized}}
\def\Complexesstack{\mathcal{C}\!\mathit{omplexes}}
% \Pic is the operator that assigns to X its picard group, usage \Pic(X)
% \Picardstack_{X/B} denotes the Picard stack of X over B
% \Picardfunctor_{X/B} denotes the Picard functor of X over B
\def\Pic{\mathop{\mathrm{Pic}}\nolimits}
\def\Picardstack{\mathcal{P}\!\mathit{ic}}
\def\Picardfunctor{\mathrm{Pic}}
\def\Deformationcategory{\mathcal{D}\!\mathit{ef}}

% --- appended 2026-10-12 (Chief): the project own definitions of the reference/slogan/history
% display environments, copied mechanically from the shared Stacks preamble lines 8-12
% (\usepackage{verbatim} plus the three \newenvironment lines, which the chapter preamble keeps
% verbatim from the project preamble). Nothing added, removed or redefined: these environments
% carry no mathematical content and the project itself typesets them as comments. Without these
% three lines a counted statement whose own text uses them fails with Environment slogan undefined.
\usepackage{verbatim}
\newenvironment{reference}{\comment}{\endcomment}
\newenvironment{slogan}{\comment}{\endcomment}
\newenvironment{history}{\comment}{\endcomment}
\usepackage{cleveref}
\crefname{originallemma}{original lemma}{original lemmas}
\Crefname{originallemma}{Original lemma}{Original lemmas}
\makeatletter
\newlabel{more-algebra-lemma-blowup-regular-sequence}{{0BIQ}{source}{}{sourcepointer-1}{}}
\newlabel{more-algebra-lemma-blowup-regular-sequence@cref}{{[originallemma][0][]0BIQ}{[1][0][]source}{}{}{}}
\makeatother
\crefname{originallemma}{original lemma}{original lemmas}
\Crefname{originallemma}{Original lemma}{Original lemmas}
\makeatletter
\newlabel{lemma-regular-koszul-regular}{{062F}{source}{}{sourcepointer-2}{}}
\newlabel{lemma-regular-koszul-regular@cref}{{[originallemma][0][]062F}{[1][0][]source}{}{}{}}
\makeatother
\crefname{originallemma}{original lemma}{original lemmas}
\Crefname{originallemma}{Original lemma}{Original lemmas}
\makeatletter
\newlabel{lemma-koszul-regular-H1-regular}{{0CEM}{source}{}{sourcepointer-3}{}}
\newlabel{lemma-koszul-regular-H1-regular@cref}{{[originallemma][0][]0CEM}{[1][0][]source}{}{}{}}
\makeatother
\crefname{originallemma}{original lemma}{original lemmas}
\Crefname{originallemma}{Original lemma}{Original lemmas}
\makeatletter
\newlabel{algebra-lemma-affine-blowup-quotient-description}{{0G8S}{source}{}{sourcepointer-4}{}}
\newlabel{algebra-lemma-affine-blowup-quotient-description@cref}{{[originallemma][0][]0G8S}{[1][0][]source}{}{}{}}
\makeatother
\crefname{originallemma}{original lemma}{original lemmas}
\Crefname{originallemma}{Original lemma}{Original lemmas}
\makeatletter
\newlabel{lemma-functorial}{{0624}{source}{}{sourcepointer-5}{}}
\newlabel{lemma-functorial@cref}{{[originallemma][0][]0624}{[1][0][]source}{}{}{}}
\makeatother
\crefname{originallemma}{original lemma}{original lemmas}
\Crefname{originallemma}{Original lemma}{Original lemmas}
\makeatletter
\newlabel{lemma-koszul-regular-flat-base-change}{{062H}{source}{}{sourcepointer-6}{}}
\newlabel{lemma-koszul-regular-flat-base-change@cref}{{[originallemma][0][]062H}{[1][0][]source}{}{}{}}
\makeatother
\crefname{originallemma}{original lemma}{original lemmas}
\Crefname{originallemma}{Original lemma}{Original lemmas}
\makeatletter
\newlabel{lemma-cone-koszul}{{0629}{source}{}{sourcepointer-7}{}}
\newlabel{lemma-cone-koszul@cref}{{[originallemma][0][]0629}{[1][0][]source}{}{}{}}
\makeatother
\title{A presentation of the blowup algebra of an H1-regular sequence}
\author{The Stacks Project Authors}
\date{stacks-project:more-algebra.tex:local-snapshot-20261012; GNU FDL 1.2 or later}
\begin{document}
\maketitle
\setcounter{subsection}{1}
\noindent\textit{Editorial note.} All delivered mathematical text is the author's own, copied from the pinned Stacks Project chapter file more-algebra.tex; no statement or proof step was rewritten, reordered or removed.\par\medskip
\noindent\textit{Editorial note.} The chapter file pinned-original-source.tex carries only \textbackslash{}input\{preamble\} on its first line, so no preamble text lives inside the chapter. The retained file sources/stacks-more-algebra/retained-preamble.tex is a mechanical, line-for-line copy of the project's single shared preamble as staged in the snapshot (Xy-pic setup plus theorem declarations and mathematical macros). Nothing was added, removed or redefined, and the excerpt serves only as preamble, never as a span coordinate.\par\medskip
\noindent\textit{Editorial note.} Both counted spans start at the \textbackslash{}begin line of their environment, so the statement and the author's proof are delivered whole; surrounding prose of the section is deliberately outside the counted spans and is carried as standard prerequisites instead.\par\medskip
\noindent\textit{Editorial note.} The citation fields (source\_id, source\_version, authors, source\_url, license.name, source\_title) are identical to those of every other card this lane has produced from this chapter, so the batch selector treats them as one source.\par\medskip
\noindent\textit{Allowed standard prerequisites (editorial).}\begin{itemize}
\item A regular sequence is Koszul regular and a Koszul regular sequence is H1-regular. Lemmas \textbackslash{}ref\{lemma-regular-koszul-regular\} and \textbackslash{}ref\{lemma-koszul-regular-H1-regular\}, cited at proof lines 7885-7887. Explain why the hypothesis of the lemma covers the classical cases of regular and Koszul regular sequences.
\item The blowup algebra is the quotient of the polynomial ring over R by the ideal of torsion elements. Algebra, Lemma \textbackslash{}ref\{algebra-lemma-affine-blowup-quotient-description\}, cited at proof line 7894. Reduces the identification of the blowup algebra to the absence of torsion by the first element.
\item Koszul complexes are functorial under an isomorphism of the underlying modules, and Koszul regularity is preserved by flat base change. Lemma \textbackslash{}ref\{lemma-functorial\} and Lemma \textbackslash{}ref\{lemma-koszul-regular-flat-base-change\}, cited at proof lines 7915-7917. Transfer H1-regularity from the given sequence to the transformed sequence in the polynomial ring.
\item The Koszul complex of a, ay\_2 - a\_2, ..., ay\_r - a\_r is the cone on multiplication by a on the Koszul complex of the remaining elements. Lemma \textbackslash{}ref\{lemma-cone-koszul\}, cited at proof lines 7918-7919. Turns vanishing of the first homology of the Koszul complex into injectivity of multiplication by a on the zeroth homology.
\end{itemize}\medskip
\ifdefined\hypertarget\hypertarget{sourcepointer-1}{}\fi
\noindent\textit{Source reference.} Original lemma 0BIQ: focused-seeds/AGP/sources/stacks-snapshot-20261011/tags.txt, source lines 14498--14498. The lemma's own canonical tag, quoted so that readers can resolve the label inside the delivered text; it points at the counted statement itself and carries no proof content.\par\medskip
\ifdefined\hypertarget\hypertarget{sourcepointer-2}{}\fi
\noindent\textit{Source reference.} Original lemma 062F: focused-seeds/AGP/sources/stacks-snapshot-20261011/tags.txt, source lines 7794--7794. Navigation pointer to a statement the counted proof refers to but does not contain; nothing load-bearing is delegated to it.\par\medskip
\ifdefined\hypertarget\hypertarget{sourcepointer-3}{}\fi
\noindent\textit{Source reference.} Original lemma 0CEM: focused-seeds/AGP/sources/stacks-snapshot-20261011/tags.txt, source lines 15582--15582. Navigation pointer to a statement the counted proof refers to but does not contain; nothing load-bearing is delegated to it.\par\medskip
\ifdefined\hypertarget\hypertarget{sourcepointer-4}{}\fi
\noindent\textit{Source reference.} Original lemma 0G8S: focused-seeds/AGP/sources/stacks-snapshot-20261011/tags.txt, source lines 20277--20277. Navigation pointer to a statement the counted proof refers to but does not contain; nothing load-bearing is delegated to it.\par\medskip
\ifdefined\hypertarget\hypertarget{sourcepointer-5}{}\fi
\noindent\textit{Source reference.} Original lemma 0624: focused-seeds/AGP/sources/stacks-snapshot-20261011/tags.txt, source lines 7783--7783. Navigation pointer to a statement the counted proof refers to but does not contain; nothing load-bearing is delegated to it.\par\medskip
\ifdefined\hypertarget\hypertarget{sourcepointer-6}{}\fi
\noindent\textit{Source reference.} Original lemma 062H: focused-seeds/AGP/sources/stacks-snapshot-20261011/tags.txt, source lines 7796--7796. Navigation pointer to a statement the counted proof refers to but does not contain; nothing load-bearing is delegated to it.\par\medskip
\ifdefined\hypertarget\hypertarget{sourcepointer-7}{}\fi
\noindent\textit{Source reference.} Original lemma 0629: focused-seeds/AGP/sources/stacks-snapshot-20261011/tags.txt, source lines 7788--7788. Navigation pointer to a statement the counted proof refers to but does not contain; nothing load-bearing is delegated to it.\par\medskip
\begin{lemma}
\label{lemma-blowup-regular-sequence}
Let $a, a_2, \ldots, a_r$ be an $H_1$-regular sequence in a ring $R$
(for example a Koszul regular sequence or a regular sequence, see
Lemmas \ref{lemma-regular-koszul-regular} and
\ref{lemma-koszul-regular-H1-regular}).
With $I = (a, a_2, \ldots, a_r)$ the blowup algebra $R' = R[\frac{I}{a}]$
is isomorphic to $R'' = R[y_2, \ldots, y_r]/(a y_i - a_i)$.
\end{lemma}

\begin{proof}
By Algebra, Lemma \ref{algebra-lemma-affine-blowup-quotient-description}
it suffices to show that $R''$ is $a$-torsion free.

\medskip\noindent
We claim $a, ay_2 - a_2, \ldots, ay_n - a_r$ is a $H_1$-regular
sequence in $R[y_2, \ldots, y_r]$. Namely, the map
$$
(a, ay_2 - a_2, \ldots, ay_n - a_r) :
R[y_2, \ldots, y_r]^{\oplus r}
\longrightarrow
R[y_2, \ldots, y_r]
$$
used to define the Koszul complex on $a, ay_2 - a_2, \ldots, ay_n - a_r$
is isomorphic to the map
$$
(a, a_2, \ldots, a_r) :
R[y_2, \ldots, y_r]^{\oplus r} \longrightarrow
R[y_2, \ldots, y_r]
$$
used to the define the Koszul complex on $a, a_2, \ldots, a_r$
via the isomorphism
$$
R[y_2, \ldots, y_r]^{\oplus r}
\longrightarrow
R[y_2, \ldots, y_r]^{\oplus r}
$$
sending $(b_1, \ldots, b_r)$ to
$(b_1 - b_2y_2 \ldots - b_ry_r, -b_2, \ldots, - b_r)$.
By Lemma \ref{lemma-functorial} these Koszul complexes are isomorphic.
By Lemma \ref{lemma-koszul-regular-flat-base-change}
applied to the flat ring map $R \to R[y_2, \ldots, y_r]$
we conclude our claim is true. By Lemma \ref{lemma-cone-koszul}
we see that the Koszul complex $K$ on $a, ay_2 - a_2, \ldots, ay_n - a_r$
is the cone on $a : L \to L$ where $L$ is the Koszul complex
on $ay_2 - a_2, \ldots, ay_n - a_r$. Since $H_1(K) = 0$ by the claim,
we conclude that $a : H_0(L) \to H_0(L)$ is injective, in other words that
$R'' = R[y_2, \ldots, y_r]/(a y_i - a_i)$
has no nonzero $a$-torsion elements as desired.
\end{proof}

\end{document}
